\section{Quality and efficiency of the implementation}
\label{sec:implementation}

\HECriterion{Quality and efficiency of the implementation --- aspects to be
taken into account}{%
\begin{itemize}
  \item Quality and effectiveness of the work plan, assessment of risks, and
  appropriateness of the effort assigned to work packages and the resources
  overall.
  \item Capacity and role of each participant, and the extent to which the
  consortium as a whole brings together the necessary expertise.
\end{itemize}}

\subsection{Work plan and resources\HEBudget{e.g. \HEImplementationBudget,
including tables}}
\label{sec:work-plan}
\begin{heguidance}
Briefly present the overall work-plan structure. Show the timing of the work
packages and their components with a Gantt chart or equivalent, and their
interrelationships with a PERT chart or equivalent.

Provide full details based on the project's logical structure and stages.
Each work package should be a substantial part of the work plan; their number
should be proportionate to the project's scale and complexity. Break work
packages into tasks or another justified structure. Describe and quantify the
activities, explain why they are needed and who performs them, and give enough
detail to justify the requested personnel, equipment, consumables, meetings and
other resources and to allow progress monitoring.

Align resources with work-package objectives and deliverables. A distinct
project-management work package is advised. Give data management,
dissemination, exploitation and communication clear visibility through tasks
or work packages. Plan updates of the data management plan and the plan for
dissemination and exploitation including communication, recording both
completed and planned activities where these plans are required.

Provide the work packages (3.1a), their descriptions (3.1b), deliverables
(3.1c), milestones (3.1d), critical implementation risks and mitigation (3.1e),
staff effort (3.1f), subcontracting justifications (3.1g), and equipment
justifications where equipment costs exceed 15\% of a participant's personnel
costs (3.1h). Keep the resource figures consistent with the budget in Part A.
\end{heguidance}

\HEParagraph{Work-plan structure}
\HEPlaceholder{Explain the work packages, their dependencies, sequencing and
contribution to the objectives.}

\HEParagraph{Timing and interrelationships}
\HEPlaceholder{Insert a legible Gantt chart and a PERT chart or equivalent
diagram, with work-package identifiers consistent with the tables below.}

\begin{heguidance}
Use plain text in the section 3.1 tables. If invited to grant-agreement
preparation, these tables must be encoded in the grant-management IT tool,
which does not support graphics or special formats.
\end{heguidance}

\HETableTitle{3.1a}{List of work packages}
\label{tab:work-packages}
\noindent
\begin{tblr}{
  width=\linewidth,
  colspec={X[0.7,l] X[2,l] X[0.9,l] X[1.1,l] X[0.9,l] X[0.7,l] X[0.7,l]},
  hlines, vlines,
  row{1}={font=\bfseries,bg=HETableGrey},
  colsep=3bp,
}
{Work\\package\\No.} & Work package title & {Lead\\participant\\No.} &
{Lead\\participant\\short name} & {Person-\\months} & {Start\\month} &
{End\\month} \\
1 & \HEPlaceholder{title} & \HEPlaceholder{No.} & \HEPlaceholder{short name} &
\HEPlaceholder{PM} & \HEPlaceholder{M} & \HEPlaceholder{M} \\
2 & \HEPlaceholder{title} & \HEPlaceholder{No.} & \HEPlaceholder{short name} &
\HEPlaceholder{PM} & \HEPlaceholder{M} & \HEPlaceholder{M} \\
3 & \HEPlaceholder{title} & \HEPlaceholder{No.} & \HEPlaceholder{short name} &
\HEPlaceholder{PM} & \HEPlaceholder{M} & \HEPlaceholder{M} \\
\end{tblr}
\begin{heguidance}
Add or remove rows to match the project. Start and end months are measured from
the project start; month 1 is the first project month. The person-month total
for each work package must agree with table 3.1f.
\end{heguidance}


\HETableTitle{3.1b}{Work package description}
\label{tab:work-package-description}
\begin{heguidance}
Repeat the following block for each work package. Participants and effort appear
in table 3.1f; the lead participant and start/end months appear in table 3.1a.
\end{heguidance}

% Duplicate this complete block for WP2, WP3, and any further work packages.
% The description box is breakable and can grow across pages.
\noindent
\begin{tabularx}{\linewidth}{|L{45mm}|Y|}
\hline
\textbf{Work package number} & \HEPlaceholder{WP number} \\
\hline
\textbf{Work package title} & \HEPlaceholder{WP title} \\
\hline
\end{tabularx}
\begin{tcolorbox}[enhanced,breakable,colback=white,colframe=black,
  boxrule=0.4bp,arc=0pt,left=5bp,right=5bp,top=5bp,bottom=5bp,
  before skip=0pt,after skip=4bp]
\textbf{Objectives}\par
\HEPlaceholder{State the work-package objectives and their relationship to
the overall project objectives.}
\par\medskip
\textbf{Description of work}\par
\begin{heguidance}
Where appropriate, break the work down into tasks. Identify the lead partner
and the role of each participant. For each task, quantify the amount of work
and explain why it is needed and who will perform it. Provide sufficient detail
to justify the resources. List linked deliverables in table 3.1c without
repeating their details here.
\end{heguidance}
\HEParagraph{Task [WP].[task number]: [task title]}
\HEPlaceholder{Identify the lead and contributing participants, timing,
activities, quantifiable work, dependencies and resource justification.}
\end{tcolorbox}


\HETableTitle{3.1c}{List of Deliverables}
\label{tab:deliverables}
\begin{heguidance}
Include only deliverables essential for effective project monitoring. Include
a data management plan (DMP) and a plan for dissemination and exploitation
including communication activities as distinct deliverables within the first
6 months. The DMP must evolve during the project.
\end{heguidance}

\noindent
\begin{tblr}{
  width=\linewidth,
  colspec={X[0.8,l] X[1.55,l] X[1.8,l] X[0.8,l] X[1.1,l] X[0.8,l] X[1.05,l] X[0.8,l]},
  hlines, vlines,
  row{1}={font=\bfseries,bg=HETableGrey},
  colsep=2bp,
}
Number & {Deliverable\\name} & {Short\\description} & {Work\\package\\number} &
{Short name\\of lead\\participant} & Type & {Dissemi-\\nation\\level} &
{Delivery\\date\\(months)} \\
{[WP].1} & \HEPlaceholder{name} & \HEPlaceholder{description} & \HEPlaceholder{WP} &
\HEPlaceholder{name} & \HEPlaceholder{code} & \HEPlaceholder{code} & \HEPlaceholder{M} \\
{[WP].2} & \HEPlaceholder{name} & \HEPlaceholder{description} & \HEPlaceholder{WP} &
\HEPlaceholder{name} & \HEPlaceholder{code} & \HEPlaceholder{code} & \HEPlaceholder{M} \\
\end{tblr}

\begin{heguidance}
\textbf{Key.} Order deliverables by delivery date. Number them as
\texttt{<WP number>.<number within WP>}; for example, 4.2 is the second
deliverable from work package 4. Delivery dates are months from project start
(month 1).

\textbf{Type codes:} R = document or report, excluding periodic/final reports;
DEM = demonstrator, pilot, prototype or plan designs; DEC = websites, patent
filings, press/media actions or videos; DATA = data sets or microdata;
DMP = data management plan; ETHICS = ethics deliverables; SECURITY = security
deliverables; OTHER = software, technical diagrams, algorithms, models, etc.

\textbf{Dissemination levels:} PU = public and fully open (published on the
CORDIS project page); SEN = sensitive, limited under the Grant Agreement;
Classified R-UE/EU-R = EU RESTRICTED; Classified C-UE/EU-C = EU CONFIDENTIAL;
Classified S-UE/EU-S = EU SECRET, under Commission Decision 2015/444.
\end{heguidance}


\HETableTitle{3.1d}{List of milestones}
\label{tab:milestones}
\noindent
\begin{tblr}{
  width=\linewidth,
  colspec={X[0.8,l] X[1.5,l] X[1.1,l] X[1,l] X[2,l]},
  hlines, vlines,
  row{1}={font=\bfseries,bg=HETableGrey},
}
{Milestone\\number} & Milestone name & {Related work\\package(s)} &
{Due date\\(month)} & Means of verification \\
MS1 & \HEPlaceholder{name} & \HEPlaceholder{WP(s)} & \HEPlaceholder{M} &
\HEPlaceholder{evidence / indicator} \\
MS2 & \HEPlaceholder{name} & \HEPlaceholder{WP(s)} & \HEPlaceholder{M} &
\HEPlaceholder{evidence / indicator} \\
\end{tblr}
\begin{heguidance}
Due dates are months from project start (month 1). Explain how attainment will
be verified and refer to indicators where appropriate; for example, a working
laboratory prototype, a software release validated by a user group, or a
completed field survey with validated data quality.
\end{heguidance}


\HETableTitle{3.1e}{Critical risks for implementation}
\label{tab:risks}
\noindent
\begin{tblr}{
  width=\linewidth,
  colspec={X[2,l] X[0.9,l] X[2,l]},
  hlines, vlines,
  row{1}={font=\bfseries,bg=HETableGrey},
}
{Description of risk; indicate (i) likelihood and (ii) severity:
Low / Medium / High} & Work package(s) involved & Proposed risk-mitigation measures \\
\HEPlaceholder{risk; likelihood; severity} & \HEPlaceholder{WP(s)} &
\HEPlaceholder{mitigation measures} \\
\HEPlaceholder{risk; likelihood; severity} & \HEPlaceholder{WP(s)} &
\HEPlaceholder{mitigation measures} \\
\end{tblr}
\begin{heguidance}
A critical risk is a plausible event or issue with a potentially high adverse
impact on the project's ability to achieve its objectives. Likelihood is the
estimated probability of occurrence after accounting for the mitigating
measures. Severity expresses the relative seriousness and significance of the
effect. Rate each as low, medium or high. The list and mitigation measures can
be updated as the project progresses.
\end{heguidance}


% Keep the heading and guidance with the initial staff-effort table.
\Needspace{16\baselineskip}
\HETableTitle{3.1f}{Summary of staff effort}
\label{tab:staff-effort}
\begin{heguidance}
Show person-months over the whole project for every participant and work
package. Identify the work-package leader by putting that participant's
person-month figure for the work package in bold. Add/remove WP columns and
participant rows as needed. Check all row/column totals against table 3.1a and
the Part A budget. With \texttt{make BUDGET\_XLSX=budget.xlsx}, this table is
generated from the partner worksheets; identify WP leaders there.
\end{heguidance}
\HEStaffEffortTable


\HETableTitle{3.1g}{`Subcontracting costs' items}
\label{tab:subcontracting}
\begin{heguidance}
For each participant, describe and justify the tasks to be subcontracted.
Core project tasks should not be subcontracted. Repeat this block for each
relevant participant, or state that no subcontracting is planned.
\end{heguidance}
\noindent
\begin{tblr}{
  width=\linewidth,
  colspec={X[1.2,l] X[0.7,r] X[3,l]},
  hlines, vlines,
  row{1,2}={font=\bfseries},
  row{2}={bg=HETableGrey},
}
\SetCell[c=3]{l} \HEPlaceholder{participant number / short name} & & \\
& Cost (\texteuro) & Description of tasks and justification \\
Subcontracting & \HEPlaceholder{EUR} & \HEPlaceholder{tasks and justification} \\
\end{tblr}


\HETableTitle{3.1h}{`Purchase costs' items (major equipment costs)}
\label{tab:equipment}
\begin{heguidance}
Provide details of the need for equipment costs when they exceed 15\% of the
personnel costs of that participant, according to the budget in Part A.
Repeat this block for each relevant participant, or state that no participant
exceeds the threshold. Template V5.1 requests major equipment here.
\end{heguidance}
\noindent
\begin{tblr}{
  width=\linewidth,
  colspec={X[1.2,l] X[0.7,r] X[3,l]},
  hlines, vlines,
  row{1,2}={font=\bfseries},
  row{2}={bg=HETableGrey},
}
\SetCell[c=3]{l} \HEPlaceholder{participant number / short name} & & \\
& Cost (\texteuro) & Justification \\
Equipment & \HEPlaceholder{EUR} & \HEPlaceholder{need for equipment} \\
\end{tblr}



\subsection{Capacity of participants and consortium as a whole\HEBudget{e.g.
3 pages}}
\label{sec:consortium}
\begin{heguidance}
Individual participants are described in Part A; focus here on the consortium
as a whole and the participants' contributions to this project.
\begin{itemize}
  \item Describe how the consortium matches the objectives and brings
  together the necessary disciplinary and interdisciplinary knowledge,
  including social sciences and humanities, open science and gender expertise
  where appropriate. Include affiliated entities and associated partners.
  \item Explain access to the critical infrastructure required for the work.
  \item Explain how members complement one another and cover the value chain
  where appropriate.
  \item Show that each participant has a valid role and adequate resources
  to fulfil it.
  \item Where applicable, describe industrial or commercial involvement
  and how it supports the exploitation measures in section 2.2.
  \item If a participant requesting EU funding is based in a country, or is
  an international organisation, that is not automatically eligible, explain
  why its participation is essential to successful project implementation.
  Check eligibility against the applicable Work Programme General Annexes B.
\end{itemize}
\end{heguidance}

\HEParagraph{Complementarity and contributions}
\HEPlaceholder{Explain the consortium's collective expertise, participant
roles, complementarity, resources and access to infrastructure.}

\HEParagraph{Exploitation capacity and essential participation}
\HEPlaceholder{Describe the relevant industrial or commercial involvement
and, where applicable, justify funding for essential participants not
automatically eligible.}

